\documentclass[10pt,a4paper]{amsart}
\usepackage[margin=19mm]{geometry}
\usepackage{amsmath,amssymb,mathtools}
\usepackage[scr=rsfs,cal=euler]{mathalfa}
\usepackage{xfrac}
\usepackage[unicode,hidelinks,bookmarks=false]{hyperref}
\newcommand{\ie}{i.e.\ }
\newcommand{\cf}{cf.\ }
\newcommand{\resp}{respectively\ }
\newcommand{\Subsection}{\subsection}
\providecommand{\nobreakdash}{\nobreak\hbox{-}}
\setlength{\emergencystretch}{2em}
\multlinegap0pt
\usepackage{amssymb}
\usepackage{float}
\usepackage{xcolor}
\usepackage{tikz}
\usepackage{enumerate}
\usepackage{bm}
\newtheorem{lemma}{Lemma}
\title{A maximum matching-based refinement of Brouwer’s conjecture}
\author{Tahir Shamsher}
\date{Source: arXiv:2609.12673v1 (CC BY 4.0)}
\begin{document}
\maketitle

\noindent\emph{Source and licence.} A maximum matching-based refinement of Brouwer’s conjecture. Version arXiv:2609.12673v1 (CC BY 4.0), distributed by arXiv under the Creative Commons Attribution 4.0 International licence, http://creativecommons.org/licenses/by/4.0/, as recorded in the arXiv licence metadata for this exact version and shown on the abstract page https://arxiv.org/abs/2609.12673v1. Full licence notice: \texttt{licenses/arxiv-2609.12673-CC-BY-4.0.txt}.\par\smallskip

\noindent\textit{Editorial scope.} Counted unit: the article's lemma brotau (source0.tex lines 563--564). The article's complete original proof of it is delivered with it (source0.tex lines 565--580, one \texttt{proof} environment of 16 lines). No mathematics is written, repaired or re-derived: the statement and the proof are the article's own lines, in the article's own notation, and no retained line is substituted or re-flowed.  Retained uncounted as the source-local context the proof needs: lines 400--402; lines 532--541.  Nothing was omitted from any retained span, so the package carries no omission record.  No retained line contains a citation, so the wrapper carries no bibliography and no bibliography entry.  No retained line contains a figure environment, an image-inclusion command or drawing code, so no image is delivered and the package declares no asset dependency.  Editorial formatting only, in addition: an AMS \texttt{amsart} preamble that restates the article's own declarations and macros used by the retained lines, namely the environments \texttt{lemma} on separate counters, together with the standard packages those lines need.  The retained environments print in this document as Lemma 1 in order of appearance, whereas the article prints the same items with the numbering of its own full document.  The complete native source of the article as distributed in the arXiv e-print for this exact version is bundled unchanged as \texttt{sources/210137/source0.tex}.
		\begin{equation}\label{tauineq}
			\mathcal{S}_k(G)\leq \sum_{i=1}^{\tau(G)}\mathcal{S}_k(G_i)=\sum_{i=1}^{\tau(G)} (\Delta(G_i) + k)=e(G)+k\tau(G).
		\end{equation} 

\begin{lemma} \label{revchen} 
Let $p$ be a positive integer with $1\leq p\leq n-2$. If the Brouwer's conjecture holds for a graph $G$ when $k=n-p-1$, that is, for the graph $G$ with $n$ vertices and $e(G)$ edges,
\begin{equation*}\label{an-p-1}
\mathcal{S}_{n-p-1}(G)\leq e(G)+\binom{n-p}{2}, 
\end{equation*}
then the Brouwer's conjecture also holds for $k=p$, that is, for the same graph $G$,
\begin{equation*}\label{p} 	
\mathcal{S}_{p}(G) \leq e(G)+\binom{p + 1}{2}.
\end{equation*}
\end{lemma}

\begin{lemma}\label{brotau} Let $G$ be a graph with $n$ vertices and $e(G)$ edges having covering number $\tau(G)$. For a positive integer $p$, if $n\geq 2\tau(G)+p$, then $$\mathcal{S}_k(G)\leq e(G)+\binom{k+1}{2}$$ holds for all integers $k$ satisfying $1\leq k\leq p$.
\end{lemma}

\begin{proof}  

Using Equation (\ref{tauineq}), we get
\begin{equation*}\label{xtauineq}
\mathcal{S}_k(G)\leq \sum_{i=1}^{\tau(G)}\mathcal{S}_k(G_i)=\sum_{i=1}^{\tau(G)} (\Delta(G_i) + k)=e(G)+k\tau(G).
\end{equation*} 
Now, for $k=n-\alpha-1 $ $(1\leq \alpha \leq p)$, assume that
\begin{equation}\label{maineq}
\mathcal{S}_k(G)\leq e(G)+k\tau(G) \leq e(G)+\binom{k+1}{2}.
\end{equation} 
To ensure the last inequality holds in \ref{maineq}, we require
$$(n-\alpha-1)\tau(G)\leq \binom{n-\alpha}{2},$$  
which simplifies to
$$n\geq 2\tau(G)+\alpha.$$  
Thus, under the condition $n\geq 2\tau(G)+p $, we have  $n\geq 2\tau(G)+p\geq 2\tau(G)+\alpha$,  therefore the result follows by Lemma \ref{revchen}. This completes the proof. 	  
\end{proof}

\end{document}
